\subsection*{7.3 Application: Evaluation of Lebesgue--Stieltjes Integrals}

As an application of the dominated convergence theorem, we provide a sufficient condition under which a Lebesgue integral on the real number line can be computed as a Riemann--Stieltjes integral. This connection allows us to evaluate abstract Lebesgue--Stieltjes (LS) integral using tools from calculus. Throughout the rest of this section, we denote the Riemann--Stieltjes (RS) integral of a function $g(x)$ with respect to a Stieltjes measure function $F(x)$ from $a$ to $b$ by $\int_a^b g(x)\, dF(x)$.

The next theorem is a link between the Lebesgue--Stieltjes integral and the Riemann--Stieltjes integral.

\begin{thmbox}{7.8}
Consider a measure space $(\mathbb{R},\mathcal{B}(\mathbb{R}),\mu)$ associated with the Stieltjes measure function $F(x)=\mu((-\infty,x])$. Assume that $\mu(\{a\})=0$, i.e., assume $F(x)$ is continuous at $x=a$. If $g$ is a continuous function on $[a,b]$, then $g$ is $\mu$-integrable on $[a,b]$ and we have
\[
\int_{[a,b]} g\, d\mu = \int_a^b g(x)\, dF(x),
\]
where the left-hand side is an LS integral and the right-hand side is an RS integral.
\end{thmbox}

\textit{Proof} Since $g$ is continuous on $[a,b]$, it is measurable. Furthermore, $g$ is bounded on $[a,b]$, i.e., there exists a constant $K$ such that $|g(x)|\le K$ for all $x\in [a,b]$. Therefore, the function $g$ is $\mu$-integrable because
\[
\int_{[a,b]} g^+\, d\mu \le \int_{[a,b]} K\, d\mu = K\mu([a,b])<\infty
\]
\[
\int_{[a,b]} g^-\, d\mu \le \int_{[a,b]} K\, d\mu = K\mu([a,b])<\infty.
\]

By Theorem 1.1, the Riemann--Stieltjes on the right-hand side exists since $g$ is continuous.

The idea of proof is to approximate the function $g$ by simple functions. We partition $[a,b]$ into mutually disjoint intervals
\[
\{a\},\ (a_1,b_1],\ (a_2,b_2],\dots,(a_n,b_n],
\]
with $a=a_1<b_1=a_2<b_2=\cdots a_n<b_n=b$, and define two simple functions
\[
\underline{g}(x)\triangleq g(a)1_{\{a\}}+\sum_{i=1}^{n} m_i 1_{(a_i,b_i]},
\qquad
\overline{g}(x)\triangleq g(a)1_{\{a\}}+\sum_{i=1}^{n} M_i 1_{(a_i,b_i]},
\]
where
\[
m_i \triangleq \inf_{a_i<x\le b_i} g(x)
\qquad \text{and} \qquad
M_i \triangleq \sup_{a_i<x\le b_i} g(x)
\]
for $i=1,2,\dots,n$. It is easily seen that
\[
\underline{g}\le g \le \overline{g}.
\]

Compute the Lebesgue integrals of $\underline{g}$, $g$, and $\overline{g}$ and apply the monotonic property,
\[
\sum_{i=1}^{n} m_i \mu((a_i,b_i])
=
\int_{[a,b]} \underline{g}
\le
\int_{[a,b]} g
\le
\int_{[a,b]} \overline{g}
\le
\sum_{i=1}^{n} M_i \mu((a_i,b_i]).
\]

We have used the assumption that $\mu(\{a\})=0$ in this step.

Because the Riemann--Stieltjes integral $\int_a^b g(x)\, dF(x)$ exists, the left-most and right-most sums above converge to a common limit. By a sandwich argument, the integral $\int_{[a,b]} g\, d\mu$ in the middle is equal to the Riemann--Stieltjes integral $\int_a^b g(x)\, dF(x)$. \hfill $\square$

\noindent$\blacktriangleright$ \textbf{Practical Calculations of Lebesgue--Stieltjes Integrals} Theorem 7.8 assumes that $g$ is continuous, but this assumption can be relaxed. With a little more work, one can show that:
\begin{enumerate}[label=(\alph*)]
    \item A function $g$ is Riemann--Stieltjes integrable if and only if $g$ is continuous almost everywhere with respect to the measure $\mu$ on $[a,b]$.
    \item If $g$ is Riemann--Stieltjes integrable on $[a,b]$, then $g$ is integrable with respect to the completion of the measure $\mu$ induced by the Stieltjes measure function $F$, and the two integrals are equal to each other.
\end{enumerate}

Nonetheless, Theorem 7.8 is sufficient for many practical calculations. If the function $g$ to be integrated is piece-wise continuous, we can compute the integral by dividing the $x$-axis into multiple pieces so that $g$ is continuous in each part. As long as the discontinuity points of $g(x)$ and the discontinuity points of the Stieltjes measure function $F(x)$ do not overlap, we can compute the integral of $g(x)$.

We note that the result in Theorem 7.8 also covers Lebesgue integration with respect to the Lebesgue measure $\lambda$. The Lebesgue integral $\int_{[a,b]} g\, d\lambda$ can be evaluated using the Riemann integral $\int_a^b g(x)\, dx$.

\textbf{Example 7.3.1 (Computation of Lebesgue--Stieltjes Integral)} \\
Consider a probability space $(\mathbb{R},\mathcal{B}(\mathbb{R}),P)$, where $P$ is the probability measure generated by the Stieltjes measure function
\[
F(x)=
\begin{cases}
0 & \text{if } x\le 0,\\
x^2 & \text{if } 0<x\le 1,\\
1 & \text{if } 1\le x.
\end{cases}
\]

The Lebesgue integral of a continuous function $g(x)$ defined on interval $[0,1]$ can be evaluated by
\[
\int_{[0,1]} g(x)\, dP(x)=\int_{0}^{1} g(x)\, d(x^2).
\]

When $g(x)=x$, it is the expectation of the probability distribution
\[
\int_{0}^{1} x\, d(x^2)=\int_{0}^{1} 2x^2\, dx = \frac{2}{3}[x^3]_0^1 = \frac{2}{3}.
\]

If we want to compute $\int \sin(\pi x)\, dP(x)$, we can first reduce it to a Riemann integral
\[
\int \sin(\pi x)\, dP(x)=\int_{0}^{1} \sin(\pi x)\, d(x^2)=2\int_{0}^{1} x\sin(\pi x)\, dx.
\]

Using integration by parts, we obtain the answer $2/\pi$.

\begin{thmbox}{7.9}
We continue with the notation in Theorem 7.8. Suppose $g$ is a piece-wise continuous function so that the Riemann--Stieltjes integral
\[
\int_a^b g(x)\, dF(x)
\]
exists for all $a<b$.

If one of the integrals
\begin{equation}
\lim_{a\to-\infty,\ b\to\infty} \int_a^b |g(x)|\, dF(x)
\qquad \text{or} \qquad
\int_{\mathbb{R}} |g(x)|\, d\mu(x)
\tag{7.4}
\end{equation}
exists and is finite, then the other integral exists and is finite. In this case, we have
\begin{equation}
\int_{\mathbb{R}} g(x)\, d\mu(x)=\int_{-\infty}^{\infty} g(x)\, dF(x).
\tag{7.5}
\end{equation}
\end{thmbox}

\textit{Proof} We first suppose that $\int_{-\infty}^{\infty} g(x)\, dF(x)$ exists and prove that $g(x)$ is $\mu$-integrable. We truncate the function $g(x)$ to $g(x)\cdot 1_{[-n,n]}(x)$, whose value is zero when $|x|>n$. By the monotone convergence theorem,
\[
\int_{\mathbb{R}} |g(x)|\, d\mu
=
\lim_{n\to\infty} \int |g(x)|1_{[-n,n]}\, d\mu
=
\lim_{n\to\infty} \int_{-n}^{n} |g(x)|\, dF(x)
<\infty.
\]

We have applied Theorem 7.8 in the second step above. Therefore $g\in L^1(\mu)$.

Conversely, suppose that $g(x)$ is $\mu$-integrable. This is equivalent to $\int_{\mathbb{R}} |g|\, d\mu < \infty$. Therefore, for any finite $a<b$,
\[
\int_a^b |g(x)|\, dF(x)
=
\int_{[a,b]} |g(x)|\, d\mu(x)
\le
\int_{\mathbb{R}} |g(x)|\, d\mu(x).
\]

Since $g$ is in $L^1(\mu)$, the Lebesgue integral $\int_{\mathbb{R}} |g|$ is finite. So $\int_a^b |g(x)|\, dF(x)$ is upper bounded by a constant, for all $a$ and for all $b$. Taking limit, the integral $\int_{-\infty}^{\infty} |g(x)|\, dF(x)$ is also upper bounded by the same constant and hence is finite.

Next, assume that one (and hence both) of the integrals in (7.4) is finite. Apply DCT with $X_n \triangleq g\cdot 1_{[-n,n]}$ and $Y=|g|$. We check that $|X_n|\le Y$ and $X_n\to g$ as $n\to\infty$. Since $Y$ is integrable by assumption, we can apply DCT to obtain
\[
\int_{\mathbb{R}} g(x)\, d\mu(x)
=
\lim_{n\to\infty} \int_{\mathbb{R}} g(x)\cdot 1_{[-n,n]}(x)\, d\mu(x).
\]

\[
= \lim_{n\to\infty} \int_{-n}^{n} g(x)\, dF(x)
\triangleq \int_{-\infty}^{\infty} g(x)\, dF(x).
\]

This proves that the LS integral and the improper RS integral in (7.5) have the same value. \hfill $\square$

We will demonstrate the application of the previous theorem by using it to calculate the mean of the exponential distribution, which has infinite support.

\textbf{Example 7.3.2 (Expected Value of the Exponential Distribution)} \\
The probability density function of the exponential distribution is given by
\[
f(x)\triangleq \lambda e^{-\lambda x}\cdot 1_{[0,\infty)}(x),
\]
where $\lambda$ is a positive constant. We can create an LS measure on $\mathbb{R}$ from the Stieltjes measure function
\[
F(x)=
\begin{cases}
\int_{0}^{x} \lambda e^{-\lambda t}\, dt & \text{if } x\ge 0,\\
0 & \text{if } x<0.
\end{cases}
\]

Denote the resulting LS measure by $P$. Suppose we want to compute the mean $\int_{\mathbb{R}} x\, dP(x)$. Since the negative real number line has measure zero with respect to measure $P$, we may focus on the positive real number line. To apply Theorem 7.9, we check that the RS integral
\[
\int_{a}^{b} x\, dF(x)=\int_{a}^{b} x\lambda e^{-\lambda x}\, dx
\]
exists for all $0\le a<b$, and moreover, the improper integral $\int_{0}^{\infty} x\lambda e^{-\lambda x}\, dx$ exists and is equal to $1/\lambda$, which is a finite value. Therefore $\int_{[0,\infty)} x\, dP(x)=1/\lambda$.
